\documentclass{article}

\usepackage{amsmath,amssymb}
\usepackage{xurl}
\usepackage[hidelinks]{hyperref}
\setlength{\emergencystretch}{2em}

% Shared commands for the notes in docs/paper-gaps/.
% Every note loads this file via % Shared commands for the notes in docs/paper-gaps/.
% Every note loads this file via % Shared commands for the notes in docs/paper-gaps/.
% Every note loads this file via \input{command}. Project-wide notation
% belongs here, never in per-note redefinitions.

% --- Project configuration: edit these three lines for your repository -----
\newcommand{\projectIssueBase}{https://github.com/OWNER/REPO/issues/}
\newcommand{\projectPrBase}{https://github.com/OWNER/REPO/pull/}
\newcommand{\projectDocBase}{https://OWNER.github.io/REPO/docs/}
% ---------------------------------------------------------------------------

\newcommand{\Lean}{\textsc{Lean}}

% A formal declaration name, upright sans, breakable at underscores.
\ExplSyntaxOn
\NewDocumentCommand{\leanid}{m}{
  \ifmmode
    \textsf{\detokenize{#1}}
  \else
    \group_begin:
    \urlstyle{sf}
    \tl_set:Nn \l_tmpa_tl {#1}
    \tl_replace_all:Nnn \l_tmpa_tl {\_} {_}
    \exp_args:NV \path \l_tmpa_tl
    \group_end:
  \fi
}
\ExplSyntaxOff

% Traceability links; use in footnotes, not in the running prose.
\newcommand{\ghissue}[1]{\href{\projectIssueBase#1}{GitHub issue \##1}}
\newcommand{\ghpr}[1]{\href{\projectPrBase#1}{PR \##1}}
\newcommand{\doclink}[1]{\href{\projectDocBase#1.html}{\path{#1}}}

% Domain notation (norms, operators, Dirac kets, ...) for your project's
% mathematics goes below, using \providecommand so a note-local refinement
% stays possible.
\providecommand{\tr}{\operatorname{tr}}

% Verdict marker: \gapnote{<kind>}{<status>}. Kind is the policy
% classification (clarification, local-correction, scope-restriction,
% unfaithful, false-source, open-gap); status is open, wip (elimination
% underway), resolved, or historical. Machine-read by the site index;
% prints a verdict line.
\newcommand{\gapnote}[2]{\par\noindent\textbf{Verdict:} #1 (#2).\par}
. Project-wide notation
% belongs here, never in per-note redefinitions.

% --- Project configuration: edit these three lines for your repository -----
\newcommand{\projectIssueBase}{https://github.com/OWNER/REPO/issues/}
\newcommand{\projectPrBase}{https://github.com/OWNER/REPO/pull/}
\newcommand{\projectDocBase}{https://OWNER.github.io/REPO/docs/}
% ---------------------------------------------------------------------------

\newcommand{\Lean}{\textsc{Lean}}

% A formal declaration name, upright sans, breakable at underscores.
\ExplSyntaxOn
\NewDocumentCommand{\leanid}{m}{
  \ifmmode
    \textsf{\detokenize{#1}}
  \else
    \group_begin:
    \urlstyle{sf}
    \tl_set:Nn \l_tmpa_tl {#1}
    \tl_replace_all:Nnn \l_tmpa_tl {\_} {_}
    \exp_args:NV \path \l_tmpa_tl
    \group_end:
  \fi
}
\ExplSyntaxOff

% Traceability links; use in footnotes, not in the running prose.
\newcommand{\ghissue}[1]{\href{\projectIssueBase#1}{GitHub issue \##1}}
\newcommand{\ghpr}[1]{\href{\projectPrBase#1}{PR \##1}}
\newcommand{\doclink}[1]{\href{\projectDocBase#1.html}{\path{#1}}}

% Domain notation (norms, operators, Dirac kets, ...) for your project's
% mathematics goes below, using \providecommand so a note-local refinement
% stays possible.
\providecommand{\tr}{\operatorname{tr}}

% Verdict marker: \gapnote{<kind>}{<status>}. Kind is the policy
% classification (clarification, local-correction, scope-restriction,
% unfaithful, false-source, open-gap); status is open, wip (elimination
% underway), resolved, or historical. Machine-read by the site index;
% prints a verdict line.
\newcommand{\gapnote}[2]{\par\noindent\textbf{Verdict:} #1 (#2).\par}
. Project-wide notation
% belongs here, never in per-note redefinitions.

% --- Project configuration: edit these three lines for your repository -----
\newcommand{\projectIssueBase}{https://github.com/OWNER/REPO/issues/}
\newcommand{\projectPrBase}{https://github.com/OWNER/REPO/pull/}
\newcommand{\projectDocBase}{https://OWNER.github.io/REPO/docs/}
% ---------------------------------------------------------------------------

\newcommand{\Lean}{\textsc{Lean}}

% A formal declaration name, upright sans, breakable at underscores.
\ExplSyntaxOn
\NewDocumentCommand{\leanid}{m}{
  \ifmmode
    \textsf{\detokenize{#1}}
  \else
    \group_begin:
    \urlstyle{sf}
    \tl_set:Nn \l_tmpa_tl {#1}
    \tl_replace_all:Nnn \l_tmpa_tl {\_} {_}
    \exp_args:NV \path \l_tmpa_tl
    \group_end:
  \fi
}
\ExplSyntaxOff

% Traceability links; use in footnotes, not in the running prose.
\newcommand{\ghissue}[1]{\href{\projectIssueBase#1}{GitHub issue \##1}}
\newcommand{\ghpr}[1]{\href{\projectPrBase#1}{PR \##1}}
\newcommand{\doclink}[1]{\href{\projectDocBase#1.html}{\path{#1}}}

% Domain notation (norms, operators, Dirac kets, ...) for your project's
% mathematics goes below, using \providecommand so a note-local refinement
% stays possible.
\providecommand{\tr}{\operatorname{tr}}

% Verdict marker: \gapnote{<kind>}{<status>}. Kind is the policy
% classification (clarification, local-correction, scope-restriction,
% unfaithful, false-source, open-gap); status is open, wip (elimination
% underway), resolved, or historical. Machine-read by the site index;
% prints a verdict line.
\newcommand{\gapnote}[2]{\par\noindent\textbf{Verdict:} #1 (#2).\par}


\title{Policy for Paper-Gap Notes}
\author{}
\date{2026-05-12}

\begin{document}
\maketitle

\section{Purpose}

A paper-gap note records a mathematical discrepancy between a cited source and
the statement, proof, or route used in the formal development. The source
itself should remain intact. If a proof seems to require a hypothesis not
present in the article, a different scalar, a different normalization, or a
replacement theorem from the formal library, the change should be explained in a
separate note under \path{docs/paper-gaps/}.

The note is a mathematical document first. It should be readable by a
mathematician or mathematical physicist who has not read the corresponding
issue or pull request. Traceability data such as formal declaration names,
source-file paths, GitHub issues, and line ranges should support the exposition,
usually in footnotes, rather than govern the exposition.

\section{When a note is needed}

A note is needed when the formal statement is not literally the cited
statement. The common case is an implication in the paper,
\begin{align}
    \mathcal H \Longrightarrow C,
    \tag{\path{source}}
\end{align}
while the formal proof establishes only
\begin{align}
    \mathcal H\wedge\mathcal K \Longrightarrow C.
    \tag{\path{formal}}
\end{align}
The condition $\mathcal K$ must then be accounted for. Either it is derived
from $\mathcal H$, it is a harmless local correction, or the formal theorem is
a restricted theorem and should be marked as such.

A note is also needed when the proof is mathematically different even though
the final conclusion is the same. For example, a source proof might identify a
coefficient by a finite power-sum identity,
\begin{align}
    \sum_{q=1}^{r_a}\mu_q^N
    =\sum_{q=1}^{r_b}\nu_q^N
    \qquad (1\leq N\leq m),
\end{align}
whereas the formal route might use convergence of a sequence. These are
different mechanisms, and the note should say which hypotheses each mechanism
uses.

\section{Form of the note}

The note should introduce the relevant notation before using it. It should
state the cited assertion in the special case under discussion, then display the
formal or corrected statement, and finally compare the hypotheses and
conclusions. Long quotations are usually unnecessary; a precise citation and
the mathematical formula are better.

Displayed equations should carry the calculation. Use \verb|align| for
displayed formulae in these notes, especially for multi-line arguments:
\begin{align}
    V^{(N)}(A) &= c_N V^{(N)}(B),\\
    \langle V^{(N)}(B_k),V^{(N)}(A_j)\rangle
               &\xrightarrow[N\to\infty]{}0,\\
    \dim\operatorname{span}
    \left(\{V^{(N)}(A_j)\}_j\cup\{V^{(N)}(B_k)\}\right)
               &= r_A+1.
\end{align}
The contradiction should then be explained in prose, for instance by saying
which nonzero coefficient is forced to vanish.
Source tags belong on the right-hand side of the displayed line, using
\verb|\tag| when a tag is needed. Short definitions may remain inline when the
sentence is clearer that way.

Formal identifiers should be written with \verb|\leanid| and should usually
appear in footnotes. For example, the prose may say that the formal theorem is
the same-dimensional gauge-phase component of a source lemma, and the footnote
may identify the declaration.\footnote{Formal declaration:
    \leanid{Model.correctedStatement_of_extra_hypothesis}
    in \doclink{Project/Path/To/File}.}
Do not put formal identifiers in displayed equations unless the mathematical
object itself is a formal declaration.

\section{Shared commands}

All notes in this directory should load \path{docs/paper-gaps/command.tex}.
Shared notation belongs there. In particular,
\begin{align}
    \tr(X),\qquad \leanid{someDeclaration}
\end{align}
should be produced by the shared commands \verb|\tr| and \verb|\leanid| rather
than by local redefinitions in each note. This keeps the documents visually
consistent and avoids contradictory local conventions.

\section{Naming}

A note's file name has the form \path{<key>_<topic>.tex}. The key names the
source the note examines: the authors' initials followed by a two-digit year
(\path{cpsv16}, \path{pgvwc07}, \path{spwc10}), or an established short name
for a book or review (\path{wolf}, \path{rmp}, \path{mpu}). A note that
audits the formal development itself, with no single external source, takes
a reserved key for the project itself, e.g. \path{self}. The list of registered keys is maintained alongside
the index generator, which rejects an unregistered key.

The topic part states the mathematical subject, not the history of the note.
A file name is permanent once other documents cite it; version suffixes are
not used, since a note is revised in place (see Revision below) and its
history is preserved by the repository.

\section{Classification}

The conclusion of a note should give a verdict. The usual classifications are:
a notational clarification, a local correction, a scope restriction, an
unfaithful theorem, a source claim that is false as printed, or an open gap.
The classification should follow the mathematics, not the amount of formal
work required. Each note declares its classification and its status
machine-readably with \verb|\gapnote{<kind>}{<status>}| directly after
\verb|\maketitle|: the kind in the kebab-cased vocabulary above
(\path{clarification}, \path{local-correction}, \path{scope-restriction},
\path{unfaithful}, \path{false-source}, \path{open-gap}) and the status
one of \path{open}, \path{wip}, \path{resolved}, or \path{historical};
\path{wip} marks a gap whose elimination is actively underway and still
counts as live. Severity is derived from the kind, so it is never stated
separately. A resolved note keeps its classification and changes only its
status; the published index lists live notes first and dims the resolved
ones.

A local correction records a small adjustment whose corrected statement is still
available in the cited argument. A scope restriction records a theorem proved
only in a special case of the source theorem. An unfaithful theorem is more
serious: it relies on a hypothesis or intermediate result that the cited source
does not provide. Such a theorem must carry an in-source marker pointing to the
paper-gap note. A source claim is false as printed when a counterexample
satisfies the stated source hypotheses and contradicts its conclusion. This
classification concerns the cited assertion itself, rather than a mismatch
introduced by the formal statement or proof.

\section{Revision}

A note may begin provisionally, but revisions should keep it as a coherent
account of the present mathematical understanding. It should not become a
chronological log. If later work proves that $\mathcal K$ follows from
$\mathcal H$, the note should be rewritten to show that derivation. If later
work confirms that the extra hypothesis is unavoidable, the note should state
the restricted theorem explicitly.

\bibliographystyle{alpha}
\bibliography{references}

\end{document}
